\documentclass[12pt]{article}
\usepackage[T1]{fontenc}
\usepackage[utf8]{inputenc}
\usepackage[brazil]{babel}
\usepackage[margin=1in]{geometry}
\setlength{\parindent}{0pt}
\begin{document}
\section*{Tema 27}
Processo: PEDILEF 0022551-92.2008.4.01.3600/ MT\\
Situacao: Julgado\\
Ramo: DIREITO PREVIDENCIÁRIO\\
Questao: Saber se é possível aplicar tabela progressiva prevista no art. 142 da Lei n. 8.213/91, aos casos de aposentadoria por idade urbana.\\
Tese: Aplica-se a tabela progressiva prevista no art. 142 da Lei n. 8.213/91, no caso de aposentadoria por idade urbana, considerando-se como marco temporal para apuração da carência o ano em que o segurado completa a idade mínima, ainda que contadas contribuições posteriores ao ano do cumprimento do requisito etário. Vide Súmula 44 da TNU.\\
Fonte: https://www.cjf.jus.br/phpdoc/virtus/
\end{document}
